\section{Transformer for fMRI: de secuencias de ROI a representaciones de redes cerebrales}

El caso de Karan muestra «la misma arquitectura con distinta geometría de datos». La fMRI no es clasificación de imágenes: la entrada son señales de baja frecuencia que varían en el tiempo en múltiples brain regions of interest (ROI), con mucho ruido, pocas muestras y grandes diferencias entre individuos. Para transferir con éxito hay que redefinir el token, el mask objective, la cross-attention y la downstream validation.

\subsection{Interfaz de datos: secuencias ROI-by-time}

La fMRI refleja la actividad neuronal de forma indirecta mediante la blood-oxygen-level-dependent signal. Si se escriben $R$ ROI y $T$ puntos temporales como $X\in\mathbb{R}^{T\times R}$, pueden tomarse como token las ventanas temporales o los embedding de ROI. La matriz de correlación describe la functional connectivity estática, pero difícilmente representa dependencias dinámicas, no lineales y direccionales.


\lecturefigure{slide-099.jpg}{fMRI: functional connectivity, averaging y correlation maps}{CS25 V5 official deck, p.~99}


\lecturefigure{slide-100.jpg}{Brain functional networks: partición en ROI y etiquetas de red}{CS25 V5 official deck, p.~100}


\lecturefigure{slide-101.jpg}{Early ML: matrices de correlación, características diseñadas a mano y clasificadores tradicionales}{CS25 V5 official deck, p.~101}

\begin{warningbox}{La correlación entre regiones cerebrales no equivale a conexión causal}
La matriz de correlación se ve afectada por entradas comunes, el preprocesamiento y la frecuencia de muestreo. La attention o los embedding que aprende un Transformer tampoco pueden interpretarse directamente como mecanismos causales neuronales; ante todo son representaciones predictivas, y deben validarse por separado frente a intervention, reproducción entre sujetos y outcome clínico.
\end{warningbox}

\subsection{Foundation model objective: masked prediction}

El curso adopta data-driven self-supervision: se ocultan algunos ROI/time tokens y con el resto de la actividad cerebral se predicen las regiones faltantes. Si el conjunto de máscara es $M$, puede escribirse
\[
\mathcal{L}_{\text{mask}}=\sum_{(t,r)\in M}\|x_{t,r}-\hat{x}_{t,r}\|_2^2.
\]
Este objetivo obliga al modelo a aprender dependencias entre regiones y a lo largo del tiempo, en lugar de apoyarse en etiquetas de enfermedad definidas por personas; sin embargo, el beneficio del preentrenamiento depende de la cohort diversity, del scanner/site shift y del mask design.


\lecturefigure{slide-102.jpg}{Foundation models for fMRI: masked modeling y encoder transferible}{CS25 V5 official deck, p.~102}


\lecturefigure{slide-103.jpg}{Data-driven approach: aprendizaje automático de representaciones a partir de secuencias de todo el cerebro}{CS25 V5 official deck, p.~103}

\teachervoice{00:41:34--00:45:16, Karan introduce los datos ROI-by-time directamente en el Transformer y construye un self-supervised objective mediante la reconstrucción de regiones ocultas. La idea de la clase: primero buscar una tarea que aproveche datos sin etiquetar y después discutir el downstream diagnosis.}

\subsection{Cross-attention y variantes de arquitectura}

Dados un unmasked context $C$ y target region queries $Q_T$, la cross-attention es
\[
H_T=\operatorname{softmax}\!\left(\frac{Q_TK_C^{\top}}{\sqrt{d_k}}\right)V_C.
\]
La región cerebral objetivo consulta activamente al resto de la red, en lugar de simplemente promediar todas las ROI. El modelo también puede usar distintas estructuras de encoder/decoder en las dimensiones de region, time y subject; la elección debe validarse por generalización y compute, no copiando los bloques de NLP.


\lecturefigure{slide-104.jpg}{Cross-attention: predicción de la target region oculta a partir de las regiones cerebrales visibles}{CS25 V5 official deck, p.~104}


\lecturefigure{slide-105.jpg}{fMRI model architectures: encoder-decoder y conexiones de attention}{CS25 V5 official deck, p.~105}

\begin{knowledgebox}{Lectura a fondo: la arquitectura para fMRI debe seguir los ejes de los datos, no las convenciones de NLP}
La señal cerebral tiene cuatro ejes principales de variación: time, region, subject y site. Si el token representa una ventana temporal, la attention es buena para aprender dependencias a lo largo del tiempo; si el token representa una ROI, puede aprender la interaction entre redes; si ambos se aplanan directamente, la secuencia resulta muy larga y semánticamente mezclada. La cross-attention permite que la target region consulte las context regions, y el encoder-decoder puede separar la observación de la predicción, pero todas estas estructuras deben compararse con temporal convolution simple, graph model y correlation baseline. Un mask ratio demasiado bajo vuelve la reconstrucción demasiado fácil; uno demasiado alto elimina el contexto predecible. El subject embedding puede absorber las diferencias individuales, pero también puede filtrar la identidad. El split más importante es por subject/site y no por time points aleatorios; de lo contrario, la autocorrelación entre escaneos cercanos infla los resultados. El diagrama de arquitectura es solo un hypothesis space; lo que determina la evidence quality es un protocol riguroso.
\end{knowledgebox}

\subsection{Resultados: primero la representation, después la downstream task}

Las diapositivas de resultados muestran a la vez métricas de reconstruction/modeling, network dependence y downstream classification. Al leer las figuras, primero hay que verificar si el train/test split cruza subject/site, y luego comparar correlation baseline, linear model y foundation representation; si solo se divide aleatoriamente dentro del mismo sitio, el modelo puede aprovechar shortcuts del scanner o demográficos.


\lecturefigure{slide-106.jpg}{Modeling results I: métricas de reconstrucción y agrupación por redes cerebrales}{CS25 V5 official deck, p.~106}


\lecturefigure{slide-107.jpg}{Modeling results II: comparación entre arquitecturas y regiones}{CS25 V5 official deck, p.~107}


\lecturefigure{slide-108.jpg}{Network dependence: dependencias predictivas entre distintas redes cerebrales}{CS25 V5 official deck, p.~108}


\lecturefigure{slide-109.jpg}{Downstream task: uso de embedding para caracterizar enfermedades o individuos}{CS25 V5 official deck, p.~109}


\lecturefigure{slide-110.jpg}{Downstream results: foundation representation frente a baseline tradicional}{CS25 V5 official deck, p.~110}

\begin{knowledgebox}{Lectura de la figura: orden de la evidencia en las diapositivas de resultados de fMRI}
Primero hay que ver si la representation es eficaz en la held-out signal reconstruction, después la estructura de dependencias entre las distintas network y, solo al final, la downstream disease classification. Aunque la Slide 110 muestre una accuracy mayor, todavía hay que preguntar por el subject-level split, la site leakage, la class balance y la calibration. Los gráficos de pastel o las visualizaciones de embedding solo describen cómo agrupa el modelo; no demuestran causalidad clínica. La conclusión más fuerte posible es «esta representación aporta información predictiva incremental bajo este protocol», no «el modelo entiende el cerebro» ni que pueda usarse directamente para diagnóstico.
\end{knowledgebox}

\teachervoice{00:45:16--00:48:02, el ponente explica la cross-attention como la predicción de la región objetivo a partir del contexto de todo el cerebro, e informa que la learned representation supera a los correlation/linear baselines en la downstream disease task. La nota advierte: es evidencia con datos y split específicos, no una conclusión sobre despliegue clínico.}

\begin{knowledgebox}{Niveles de validación de un foundation model en neurociencia}
\begin{enumerate}
  \item si la masked reconstruction supera a un baseline simple;
  \item si la representation es estable entre subject/site;
  \item si la downstream task se reproduce en una cohort externa;
  \item si tiene calibration, equidad y utility clínica;
  \item si la explicación del modelo concuerda con los mecanismos neuronales conocidos y con la intervention.
\end{enumerate}
Superar los dos primeros niveles no implica que los tres últimos se cumplan automáticamente.
\end{knowledgebox}

\subsection{Resumen de la sección}

El caso de fMRI muestra que la transferibilidad del Transformer proviene de la token interaction, no del nombre de NLP. La tokenización por ROI/time, el masked objective, la cross-attention y la validación entre cohort determinan juntos el éxito o el fracaso; el modelo puede aprender representaciones útiles, pero ni la attention ni la correlación pueden elevarse directamente a explicaciones causales del cerebro.
